\section{Appendix: Operational Reference}
\label{sec:appendix}

This appendix is a quick-start command reference, kept in sync with
\texttt{CLAUDE.md} and \texttt{docs/SESSION\_CONTEXT.md} at the repository
root. If any command below fails, those two files (not this report) are the
first place to check for drift.

\subsection{Build}

\begin{lstlisting}
# Standard build (must be outside any Python venv)
cd /root/ws
colcon build --symlink-install
source install/setup.bash

# If venv dependencies need to be respected:
source venv/bin/activate
python3 $(which colcon) build --symlink-install \
  --cmake-args -DPython3_EXECUTABLE=$(which python3)
source install/setup.bash
\end{lstlisting}

\subsection{Run}

\paragraph{Recommended -- web GUI (manages the full stack as child processes):}
\begin{lstlisting}
ros2 run robo_reason_gui gui_node
# open http://localhost:8080
\end{lstlisting}

\paragraph{LLM dry-run (no robot, no camera):}
\begin{lstlisting}
# Terminal 1:
ros2 launch robo_reason_bringup dry_run.launch.py use_mock_llm:=true
# Terminal 2:
ros2 run robo_reason_task_interface task_interface_node
\end{lstlisting}

\paragraph{VLM dry-run (mock PNG camera):}
\begin{lstlisting}
ros2 launch robo_reason_bringup vlm_dry_run.launch.py \
  images_dir:=/path/to/mock_frames model_name:=groq/qwen3.6-27b
\end{lstlisting}

\paragraph{Camera driver (Orbbec):}
\begin{lstlisting}
./scripts/run_orbbec_registered.sh
\end{lstlisting}

\paragraph{Camera services + overlay:}
\begin{lstlisting}
ros2 launch vlm_camera_service camera_services.launch.py \
  show_overlay:=true \
  charuco_z_sign:=1.0 \
  board_in_base_x:=-0.224 \
  board_in_base_y:=-0.348 \
  board_in_base_roll:=3.14159 \
  board_in_base_yaw:=1.5708 \
  z_offset_m:=0.01
\end{lstlisting}

\paragraph{LLM mode (real robot):}
\begin{lstlisting}
export GROQ_API_KEY=gsk_...
ros2 launch robo_reason_bringup real_robot.launch.py \
  mode:=LLM \
  reasoning_method:=fhp \
  model_name:=groq/qwen3-32b \
  temperature:=0.1
\end{lstlisting}

\paragraph{VLM mode (real robot):}
\begin{lstlisting}
export GROQ_API_KEY=gsk_...
ros2 launch robo_reason_bringup real_robot.launch.py \
  mode:=VLM \
  model_name:=groq/qwen3.6-27b \
  temperature:=0.1
\end{lstlisting}

\paragraph{Task interface CLI:}
\begin{lstlisting}
ros2 run robo_reason_task_interface task_interface_node
\end{lstlisting}

\subsection{API keys}

Set \texttt{\{PROVIDER\}\_API\_KEY} environment variables (e.g.\
\texttt{GROQ\_API\_KEY}, \texttt{NEBIUS\_API\_KEY}, \texttt{ANTHROPIC\_API\_KEY})
or add them to a \texttt{.env} file in the working directory.

\subsection{ChArUco board parameters (current, authoritative)}

\begin{center}
\begin{tabular}{@{}ll@{}}
\toprule
\textbf{Parameter} & \textbf{Value} \\
\midrule
Dictionary & \texttt{DICT\_6X6\_250} \\
Layout & $3 \times 4$ squares \\
\texttt{square\_length} & 0.062\,m \\
\texttt{marker\_length} & 0.031\,m \\
\texttt{axis\_length} & 0.08\,m \\
\bottomrule
\end{tabular}
\end{center}

\subsection{Board pose in \texttt{base\_link} (current defaults)}

\begin{center}
\begin{tabular}{@{}ll@{}}
\toprule
\textbf{Parameter} & \textbf{Value} \\
\midrule
\texttt{board\_in\_base\_x} & $-0.224$\,m \\
\texttt{board\_in\_base\_y} & $-0.348$\,m \\
\texttt{board\_in\_base\_z} & $-0.030$\,m \\
\texttt{board\_in\_base\_roll} & 3.14159\,rad ($180^\circ$ around X) \\
\texttt{board\_in\_base\_pitch} & 0.0 \\
\texttt{board\_in\_base\_yaw} & 1.5708\,rad ($90^\circ$ around Z) \\
\bottomrule
\end{tabular}
\end{center}

Intrinsic XYZ convention: roll first, then pitch, then yaw. Measured once
per physical setup -- re-measure if the camera or board mount has moved
(Section~\ref{sec:open-issues}, issue I8).

\subsection{Docker container}

\begin{center}
\begin{tabular}{@{}ll@{}}
\toprule
\textbf{Item} & \textbf{Value} \\
\midrule
Image tag & \texttt{roboreason:ur} \\
Launch alias & \texttt{roboreason-ur} (includes \texttt{--privileged -v /dev:/dev}) \\
Workspace mount & \texttt{\textasciitilde/docker\_ws/ros2\_humble} (host) $\to$ \texttt{/root/ws} (container) \\
Setup guide & \texttt{docs/docker-setup.md} \\
\bottomrule
\end{tabular}
\end{center}

\subsection{Repository pointers}

\begin{itemize}[nosep]
  \item \textbf{Detailed session log}: \texttt{docs/SESSION\_CONTEXT.md}
    (file-by-file change tables for every session).
  \item \textbf{Raw backlog}: \texttt{docs/TODO.md}.
  \item \textbf{Grasp geometry derivation}: \texttt{docs/GRASP\_GEOMETRY\_PIPELINE.md}.
  \item \textbf{Full operator guide}: \texttt{docs/ROBOREASON\_GUIDE.md}.
  \item \textbf{Docker setup for new collaborators}: \texttt{docs/docker-setup.md}.
  \item \textbf{GUI package README}: \texttt{src/robo\_reason\_gui/README.md}.
  \item \textbf{Benchmark data, plan, and figures}: \texttt{benchmark/}.
  \item \textbf{Current branch at time of writing}: \texttt{feature/refactoring}.
  \item \textbf{Superseded/historical}: the sibling working directory
    \texttt{robotic-ai/docs} contains the pre-repository VLA-era planning
    documents (Section~\ref{sec:introduction}) -- historical reference only,
    not an accurate description of the current system.
\end{itemize}
